%==============================================================
% bridges/main.tex —— 专题三：跨篇串讲（独立册）
%==============================================================
% 为什么单独成册（2026-11-02）:
%   · 定位不同: A–R 主篇与专题一/二是**新知识**, 本册是「回头贯通」的连接
%     (README.md 原话: 不是新知识点, 是篇与篇之间的连接)。
%     这也是"为什么只拆专题三、不拆专题一/二"的答案。
%   · 会持续追加: 每学到新的衔接点就加一篇, 单册便于增量维护与单独打印。
%   · 拆分前提已验证: 各篇不含任何 \ref/\eqref, 也没有 equation/align/\tag,
%     公式全为无编号 \[ \], 因此独立编译不会出现 ?? 与编号漂移。
%   · 2026-11-03: 加了 br0_skeleton(骨架) —— 分篇依据从「话题」改成「结构动作 S1–S6」,
%     篇目不再按写作顺序排, 按骨架顺序排。规约见 README.md。
%
% 编译（在 bridges/ 下执行）:
%   latexmk -pdfxe -interaction=nonstopmode main.tex
% 必须 -pdfxe; 用 -pdf 时 ctex 会报 fontset 'mac' unavailable 并卡在交互输入。
%==============================================================
\documentclass[11pt,a4paper]{article}
\usepackage[UTF8]{ctex}
\usepackage{amsmath,amssymb,amsthm}
\usepackage{geometry}
\usepackage{booktabs}
\usepackage{array}
\newcolumntype{L}[1]{>{\raggedright\arraybackslash}p{#1}}
\usepackage{enumitem}
\usepackage{hyperref}
\usepackage{titlesec}
\usepackage{fancyhdr}
\usepackage{tikz}
\usepackage{mathtools}

\geometry{margin=2.5cm}
\hypersetup{colorlinks=true,linkcolor=blue,urlcolor=blue}

% 便利宏（2026-11-03 加，与 survey/main.tex 同名同义；原四篇仍用手写定界符，不受影响）
\newcommand{\eps}{\varepsilon}
\newcommand{\abs}[1]{\left\lvert #1\right\rvert}
\newcommand{\norm}[1]{\left\lVert #1\right\rVert}

\pagestyle{fancy}
\fancyhead[L]{现代数学物理与计算机图形学 · 专题三 跨篇串讲}
\fancyhead[R]{\thepage}
\fancyfoot[C]{}

\title{\textbf{专题三：跨篇串讲}\\[4pt]
       \large 篇与篇之间的连接 —— 读法是「回头贯通」，不是「逐篇学」}
\author{}
\date{}

\begin{document}
\maketitle
\tableofcontents
\newpage

% BR0 是纲：六节点骨架 + 覆盖矩阵 + S5/S6 正文。先读它，再读 BR1–BR4。
%==============================================================
% br0_skeleton.tex —— BR0. 骨架：六个动作
%==============================================================
% 定位（2026-11-03）: 本册的**纲**。它不是第七篇串讲，而是给 BR1–BR4 指定位置。
% 分篇依据从「话题」改成「结构动作」，理由见 README.md §0。
% 规则: 每一条陈述必须是 A ≅ B；写不成的退回主篇或 survey。
% 篇幅克制: S1–S4 的正文在 BR1–BR4/S1–S4 已有, 本篇只给指针;
%           S5、S6 原本没有独立篇目, 本篇给正文(骨架级: ≅ 清单 + 判据, 不重推导)。
%==============================================================

\section{BR0. 骨架：六个动作}
\label{sec:br0}

\begin{center}
\small
\begin{tabular}{L{2.4cm}L{11.4cm}}
\toprule
\textbf{本篇 $\cong$} & \textbf{一句话} \\
\midrule
表示 $\cong$ 换坐标 & 对象不变，数组变（S1）\\
线性化 $\cong$ 取切空间 & 在一点假装线性（S2）\\
谱 $\cong$ 对角化 & 选一个让算子不耦合的基（S3）\\
流 $\cong$ 单参数群 & 走一步与走到底之差（S4）\\
约束 $\cong$ 投影 & 最近点就是正交投影（S5）\\
离散 $\cong$ 子空间逼近 & 连续对象换成有限自由度再放回去（S6）\\
\bottomrule
\end{tabular}
\end{center}

\subsection{BR0.1 为什么按动作分、不按话题分}

旧版四篇是按\textbf{话题}分的（旋转、傅里叶、PDE、概率）。按话题分有一个结构性缺陷：
\textbf{话题是开集}，每来一个新领域就要新增一篇，而新篇与旧篇之间没有咬合 —— 于是「散」不是写作问题，是分篇依据的问题。

结构动作是\textbf{闭集}：任何领域里的数学，只要被写下来，做的不外乎换坐标、取线性化、对角化、沿流走、投影、取极限这六件事之一。所以按动作分篇，篇数有上限，而新领域只能落进已有节点，或者证明节点不够。

\textbf{判据}：新增一篇的正当理由是「某个真问题归不进六个节点」；不是「这个话题有意思」。
反例：BR3 现在横跨 S1–S6 全部 —— 它不是一篇并列串讲，是骨架的纵向走一遍。读它之前该先读本篇。

\subsection{BR0.2 六节点与各自的判据}

\begin{center}
\small
\begin{tabular}{L{2.1cm}L{2.1cm}L{5.2cm}L{4.4cm}}
\toprule
\textbf{节点} & \textbf{动作} & \textbf{一句话} & \textbf{判据（判不出就不是结构）} \\
\midrule
S1 & 表示 & 换基、换坐标：对象不变，数组变 & 换基后不变量（谱、迹、内积、体积元）逐项相等 \\
S2 & 线性化 & 在一点假装线性：曲面 $\to$ 切空间 & 残差随 $\eps$ 减半降 4 倍（$O(\eps^{2})$）\\
S3 & 谱 & 选一个让算子不耦合的基 & 截断 $N$ 项的残差按 $\abs{\lambda_N}$ 衰减 \\
S4 & 流 & 沿向量场走：走一步 $\ne$ 走到底 & $\Phi_{t+s}=\Phi_t\circ\Phi_s$ 成立；$h$ 减半误差降 $2^{p}$ \\
S5 & 约束与投影 & 允许集合不是子空间时，先投影再走 & 残差与允许方向的内积为 $0$（机器精度）\\
S6 & 离散与极限 & 有限自由度逼近连续对象 & $h$ 减半误差降 $2^{p}$，$p=$ 子空间多项式次数 \\
\bottomrule
\end{tabular}
\end{center}

\textbf{六个判据的共同点}：它们都是\emph{可算的等式或比值}，不是「理解了没有」。
这是本册与主篇分工的界线：主篇负责让判据成立，本册负责指出\textbf{同一个判据在不同领域里换了什么名字}。

\begin{center}
\small
\begin{tabular}{L{2.1cm}L{5.6cm}L{5.2cm}}
\toprule
\textbf{节点} & \textbf{已有正文} & \textbf{兼在} \\
\midrule
S1 表示 & BR3 §BR3.3（张量三视角、变换律） & BR1（$S^{3}$ 与 $SO(3)$ 双覆盖）、BR4 §BR4.5 \\
S2 线性化 & BR1 §BR1.5（$\exp$ 是线性化的逆） & BR3 §BR3.4（李代数 $=T_eG$）\\
S3 谱 & BR2 全篇 & BR3 §BR3.4 末（球谐 $=$ 不可约表示）、BR4 §BR4.4 \\
S4 流 & BR3 §BR3.1 & BR1（测地线 $=$ 单参数子群）、BR2 §BR2.8（热核 $=$ 半群）\\
S5 约束与投影 & \textbf{本篇 §BR0.3（新）} & BR4 §BR4.3（条件期望 $=L^{2}$ 投影）\\
S6 离散与极限 & \textbf{本篇 §BR0.4（新）} & BR3 §BR3.2（弱形式，一段）\\
\bottomrule
\end{tabular}
\end{center}

S5、S6 之前没有独立篇目 —— 而 \texttt{NOW.md} 本期十件里七件落在它们上面。
这就是「体系缺一块」的可核对形式：不是感觉，是空格。

\subsection{BR0.3 S5：约束与投影}

\textbf{统一陈述}：给定允许集合 $C$（通常是子空间、仿射子空间或流形），
「在 $C$ 里选最好的」$=$ 「先把不许的方向扔掉，再在 $C$ 里自由地做」。
扔掉的动作就是正交投影 $P_C$；而 $C$ 定义得对不对，只看一条：
\[
\bigl\langle\, x-P_C x,\; c\,\bigr\rangle=0\qquad \forall c\in C .
\]
\textbf{这一条判据同时是八种算法的身份证明}：

\begin{enumerate}[leftmargin=1.5em]
  \item \textbf{Hodge 分解 $\cong$ 投影}：$L^{2}=d(\Lambda^{k-1})\oplus\delta(\Lambda^{k+1})\oplus\mathcal{H}^{k}$，
        即「无旋 $+$ 无散 $+$ 调和」。不可压条件 $\nabla\cdot u=0$ 不是方程，是\textbf{选了一个子空间}；
        压力 $p$ 是把动量方程投影到无散子空间时冒出来的 Lagrange 乘子。
        判据：$\langle Pu,\nabla q\rangle=0$ 对一切 $q$ 成立。
  \item \textbf{条件期望 $\cong$ $L^{2}$ 投影}：$\mathbb{E}[X\mid\mathcal{G}]$ 是 $X$ 在 $L^{2}(\Omega,\mathcal{G},P)$ 上的正交投影。
        判据：$\langle X-\mathbb{E}[X\mid\mathcal{G}],\ \mathbf{1}_A\rangle=0$，$\forall A\in\mathcal{G}$。
        与第 1 条同构：\textbf{「已知信息」就是一个子空间}，条件化就是投影到它上面。
  \item \textbf{最小二乘 $\cong$ 投影到列空间}：残差 $r=b-Ax$ 必须与 $\mathrm{range}(A)$ 正交，
        于是得到正规方程 $A^{H}r=0$。判据：$\norm{A^{H}(b-Ax)}\le \mathrm{tol}\cdot\norm{A^{H}b}$。
        QR 与 SVD 只是算这个投影的两种稳定方式。
  \item \textbf{CG $\cong$ Krylov 子空间上的 Galerkin 投影}：第 $k$ 步在 $x_0+\mathcal{K}_k(A,r_0)$ 里取
        使 $\norm{e_k}_A$ 最小的点，等价于 $r_k\perp\mathcal{K}_k$。
        判据：$\langle r_k,r_j\rangle=0$（$j<k$）逐对成立；精确算术下 $r_n=0$。
        \textbf{所以 CG 不是「迭代技巧」，是 S5 的一个实例} —— 这条把它接进了骨架，
        也解释了 $\mathrm{NOW.md}$ 第 4 条为什么只需 \texttt{apply}：投影只要「乘一下」。
  \item \textbf{正交化 $\cong$ 逐次投影到正交补}：Gram--Schmidt 的每一步都是
        $q_j=\bigl(I-\sum_{i<j}q_iq_i^{H}\bigr)v_j$ 再归一化。
        判据：$\norm{Q^{H}Q-I}\le 8n\eps$；失去正交性（MGS 无重正交）就是这个投影算错了。
  \item \textbf{Dirichlet 边界 $\cong$ 仿射约束}：解被限制在 $\{u:u|_{\partial\Omega}=g\}$ 这个\emph{仿射}子空间里。
        「行/列清零 $+$ 对角置 1」是它的显式投影，好处是\textbf{保对称}（CG 才成立）；
        用罚因子是软约束，残差留平行分量，永远差一个偏差 —— 这就是「为什么清零而不是加大罚数」。
        判据：边界节点值等于 $g$ 到机器精度，且 $A(i,j)=A(j,i)$ 逐项精确相等。
  \item \textbf{单位化四元数 $\ne$ 投影到 $SO(3)$}：$q/\norm{q}$ 是投影到 $S^{3}$，而 $SO(3)=S^{3}/\{\pm1\}$，
        所以「插值后归一化」与 slerp 走的是两条路：前者投影，后者测地线（BR1）。
        \textbf{投影 $\ne$ 测地线}，这是本节点唯一要防的误用。
  \item \textbf{保结构格式 $\cong$ 约束在等能面上（弱）}：辛格式不保持能量点值，只保持「能量有界」，
        即轨道被约束在等能面附近的一个环壳里。判据：长时能量振荡且有界；非辛格式单调漂移。
\end{enumerate}

\textbf{一句话}：看到「约束」两个字，先问「允许集合是什么，投影是哪一个」；
回答不了，就还是在用罚因子猜。

\subsection{BR0.4 S6：离散与极限}

\textbf{统一陈述}：连续问题 $u\in V$ 变成 $u_h\in V_h\subset V$，方程不再逐点成立，只要求
「在 $V_h$ 的每个方向上都投影相等」：
\[
a(u_h,v_h)=\langle f,v_h\rangle\qquad\forall v_h\in V_h .
\]
这就是 Galerkin —— \textbf{它本身就是 S5 的一次投影}，只是投到了有限维子空间。
误差由 Céa 引理一分为二：$\norm{u-u_h}_V\le\frac{M}{\alpha}\inf_{v\in V_h}\norm{u-v}_V$
（左边是离散误差，右边是\textbf{逼近误差}乘一个稳定性常数）。S6 的全部内容就是这两项。

\begin{enumerate}[leftmargin=1.5em]
  \item \textbf{弱形式 $\cong$ 内积投影}：$-u''=f$ $\Rightarrow$ $\int u'v'=\int fv$，二阶导降为一阶。
        判据：$u_h$ 的残差在 $V_h$ 上正交，即 $\langle -u_h''-f,\varphi_i\rangle=0$ 逐基函数成立。
  \item \textbf{装配 $\cong$ 同一个双线性形式换一组基}：单元刚度
        $K_e=\frac1h\begin{pmatrix}1&-1\\-1&1\end{pmatrix}$ 是 $a(\cdot,\cdot)$ 限制在 $\{\varphi_i\}$ 上的矩阵。
        判据（三条同时）：$A$ 逐项对称、未加边界时内部行和为 $0$（常数在核里）、
        1D 手算表逐项对拍 $\bigl(n=3,h=\tfrac12\Rightarrow (1/h)[[1,-1,0],[-1,2,-1],[0,-1,1]]\bigr)$。
  \item \textbf{1D P1 节点精确 $\cong$ 离散解恰好落在插值空间里}：$u=x(1-x)$、$f=2$ 时 $u_h$ 在节点上等于 $u$，
        到机器精度。这不是「精度高」，是 $u$ 二次、$f$ 常数使 $V_h$ 精确容纳了真解的那一半；
        \textbf{2D 不能套这条}（一般 $u_h\ne u$）。这是 \texttt{NOW.md} 第 3 条的「最狠判据」的数学理由。
  \item \textbf{收敛阶 $\cong$ 逼近阶}：$\norm{u-u_h}\le Ch^{p}$，$p=\min(k+1,\ s)$，
        $k$ 是单元多项式次数、$s$ 是真解的光滑度。判据：三次 $h$ 减半的误差比值落 $[3.8,4.2]$（$p=2$），
        且 $h$ 已进入渐近区 —— 否则比值乱跳，这是「阶」最常被误用的地方。
  \item \textbf{稀疏性 $\cong$ 局部性}：$\varphi_i$ 与 $\varphi_j$ 只在支撑相交时耦合，
        于是 $A$ 的非零结构 $=$ 网格的邻接图。判据：$\mathrm{nnz}=O(N)$ 而非 $O(N^{2})$（2D 三角剖分约 $7N$）。
        CSR 因此不是存储技巧，是这条局部性的必然结果 —— 它\textbf{没有数学增量}，故不进本册正文。
  \item \textbf{迭代数 $\cong$ 谱的分布}：CG 迭代数 $\lesssim\frac12\sqrt{\kappa}\log\frac2\eps$，
        而 $A$ 是 $-\Delta$ 的离散化，$\kappa(A)=O(h^{-2})$。合起来：\textbf{迭代数} $O(1/h)$。
        这一条同时解释了 \texttt{NOW.md} 第 4 条「\texttt{iters} 只记录、不设阈值」：
        它由 $\kappa$ 决定，不由实现质量决定。
  \item \textbf{采样与混叠 $\cong$ 同一件事的频域版}：截断最高频 $=$ 投影到带限子空间（$\cong$ S5）；
        采样率不够则是\textbf{投影到了错的子空间}（不是正交投影），表现为高频折叠成低频。
        判据：折叠频率（$\bmod f_s$）处的谱峰位置。
  \item \textbf{稳定性 $\cong$ 时间方向的谱半径}：显式格式要求一步算子的谱半径 $\le 1+O(\Delta t)$，
        于是 $\Delta t\le Ch^{2}$（扩散）或 $Ch$（对流，CFL）。
        判据：$\Delta t$ 越过阈值时解发散，能量单调上升而不是有界振荡。
  \item \textbf{舍入与离散的两条曲线}：$h\to0$ 时离散误差按 $h^{p}$ 降，但 $\kappa=O(h^{-2})$
        使舍入被放大，总误差在双对数图上出现\textbf{平台或谷底}。
        判据：收敛阶测试的三个点必须都在谷底左侧，否则测出来的不是 $p$。
  \item \textbf{恒等式 $\ne$ 收敛}：$d^{2}=0$、Hodge 对合 $**\omega=(-1)^{k(n-k)}\omega$、
        Stokes 用三次多项式 $+$ 张量积 Simpson 求积，都是\textbf{离散误差为零}的特例，
        所以判据可以写到 $10^{-13}$ 而不涉及 $h$。
        分清「恒等式」与「收敛」是本节点最省时间的一条经验。
\end{enumerate}

\subsection{BR0.5 工程出口覆盖（随 \texttt{NOW.md} 版本重写）}

\begin{center}
\small
\begin{tabular}{L{4.4cm}L{1.8cm}L{6.6cm}}
\toprule
\texttt{NOW.md} 条目 & 节点 & 落点 \\
\midrule
1 \texttt{grid.hpp} 结构化网格 & S6 & 本篇 §BR0.4 第 4、5 条 \\
2 \texttt{sparse.hpp} CSR & S6 & 第 5 条；\textbf{无数学增量}，只登记 \\
3 \texttt{assemble.hpp} P1 $+$ Dirichlet & S6 $+$ S5 & §BR0.4 第 1–3 条、§BR0.3 第 6 条 \\
4 \texttt{cg.hpp} & S5 $+$ S3 & §BR0.3 第 4 条、§BR0.4 第 6 条 \\
5 \texttt{npy.hpp} 落盘 & — & \textbf{无数学}，不进本册 \\
6 \texttt{math/} 四件填实现 & S3 & BR2 全篇（\texttt{fourier}、\texttt{series}、\texttt{special}、\texttt{prob}）\\
7 \texttt{Tensor} & S1 & BR3 §BR3.3 \\
8 \texttt{ode.hpp} & S4 $+$ S2 & BR3 §BR3.1（判据对照待补）\\
9 \texttt{forms.hpp} 微分形式 & S1 $+$ S5 & BR3 §BR3.3 末、§BR3.6；§BR0.4 第 10 条 \\
\bottomrule
\end{tabular}
\end{center}

\textbf{这张表的用法}：\texttt{NOW.md} 整篇重写时重跑一次，本期每件都必须有节点；
没有节点的条目要么「无数学」（写明），要么说明骨架缺一块（走 README §5 的 B4）。

\subsection{BR0.6 骨架什么时候算错}

三条自否证条件，任一触发就说明这套骨架要改，而不是这一篇要补：

\begin{enumerate}[leftmargin=1.5em]
  \item 某节点的判据凑不出来（例如 S1 说「换基不变」，却找不到一个可算的不变量）$\Rightarrow$ 该节点是装饰，删；
  \item 某个真问题（含 \texttt{NOW.md} 条目）无法归入六节点 $\Rightarrow$ 缺第七个，先登记再决定；
  \item 两个节点合并后 $\cong$ 清单不变 $\Rightarrow$ 它们本来是一个动作，合并。
\end{enumerate}

\textbf{当前状态}：S1–S4 有正文且有判据；S5、S6 在本篇有骨架级正文（$\cong$ 清单 $+$ 判据），
但\textbf{还没有一篇把它们与 BR1–BR4 横向缝起来的完整串讲} —— 这是本册下一步唯一的结构性缺口。

% 四篇彼此独立，可单抽。
%==============================================================
\section{BR1. 从四元数 slerp 到大圆测地线}
\label{sec:br1}

\begin{center}
\small
\begin{tabular}{L{3.4cm}L{10.4cm}}
\toprule
\textbf{本篇 $\cong$} & \textbf{一句话} \\
\midrule
slerp $\cong$ 球面测地线 & 等速参数化就是测地线，不是插值技巧 \\
指数映射 $\cong$ 单参数子群 & $\exp(tv)$ 是 $v$ 生成的流 \\
四元数乘 $\cong$ $S^3$ 的群结构 & 单位四元数是 $SO(3)$ 的双重覆盖 \\
大圆 $\cong$ 等距浸入 & 球面上「直」的曲线 \\
对极奇点 $\cong$ 指数映射非单射 & $q$ 与 $-q$ 是同一旋转，插值方向不唯一 \\
\bottomrule
\end{tabular}
\end{center}

\noindent\textbf{落点}：\texttt{G4}--\texttt{G6}（四元数代数 / 双覆盖 / 插值）、\texttt{D8}（辛结构：等速参数下的守恒量）
%==============================================================

本篇的目标不是讲清 slerp 怎么实现，而是把它放回几何：
\textbf{slerp 不是插值技巧，它是球面测地线的等速参数化}。
一旦看清这一点，大圆、指数映射、单参数子群、对极奇点会同时到位。

\subsection{BR1.1 问题从哪来}

两个单位四元数 $q_0,q_1\in S^3\subset\mathbb{H}$，表示两个旋转，要在其间插值。

最朴素的做法是线性插值再归一化（nlerp）：
\[
q(t)=\frac{(1-t)q_0+tq_1}{\lVert (1-t)q_0+tq_1\rVert}.
\]

它能用，但有两个病：

\begin{enumerate}
  \item \textbf{参数与角度不同步}。归一化把 $t$ 和“转过的角度”的关系弄歪了：$t$ 走一半，角度不一定走一半，角速度非均匀。
  \item \textbf{概念错位}。它把球面上的点投到弦上再投回来，等于承认“球面是弯的，但我假装它平”。
\end{enumerate}

要真正回答“球面上两点之间最直的路是什么”，必须先定义“直”。

\subsection{BR1.2 “直”在弯曲空间里的三种刻画}

球面 $S^n$ 嵌在 $\mathbb{R}^{n+1}$ 里，上面的“直线”叫\textbf{测地线}。三种等价刻画，由浅到深。

\subsubsection{（a）对称性刻画}

球面关于任何过球心的平面 $\Pi$ 反射不变。取 $\Pi=\mathrm{span}\{p,q\}$，这个反射的不动点集正是 $\Pi\cap S^n$——一个\textbf{大圆}。

\textbf{道理}：一个几何对象若被“所有保持它不变的对称”钉住，它必然最自然。这也是“为什么不是小圆”的回答：小圆不被过球心的平面反射保持。

\subsubsection{（b）变分刻画（最有用）}

设曲线 $\gamma:[0,1]\to S^n$，$\lVert\gamma\rVert=1$，扰动 $\eta\perp\gamma$（保持单位约束）。弧长
\[
L[\gamma]=\int_0^1 \lVert\dot\gamma\rVert\,dt .
\]
要求 $L$ 取驻值。分部积分（端点固定）：
\[
\delta L=\int \frac{\dot\gamma\cdot\dot\eta}{\lVert\dot\gamma\rVert}\,dt
=-\int\left(\frac{d}{dt}\frac{\dot\gamma}{\lVert\dot\gamma\rVert}\right)\cdot\eta\,dt .
\]
驻值要求括号那项垂直于所有 $\eta\perp\gamma$，即 $\dfrac{d}{dt}\dfrac{\dot\gamma}{\lVert\dot\gamma\rVert}\parallel\gamma$。

改用弧长参数（$\lVert\dot\gamma\rVert\equiv 1$）：
\[
\boxed{\ \ddot\gamma=\lambda\gamma\ }
\]
再用约束 $\lVert\gamma\rVert^2=1$ 求两次导：$\gamma\cdot\dot\gamma=0$，$\dot\gamma\cdot\dot\gamma+\gamma\cdot\ddot\gamma=0$，即 $1+\lambda=0$。于是
\[
\boxed{\ \ddot\gamma=-\gamma\ }
\]
\textbf{这是一个简谐振子方程}。解为
\[
\gamma(t)=\cos t\,p+\sin t\,v,\qquad v\perp p,\ \lVert v\rVert=1,
\]
正是平面 $\mathrm{span}\{p,v\}$ 与球面的交——大圆。半径 $R$ 的球改为 $\ddot\gamma=-\gamma/R^2$。

一般流形上的测地线方程是
\[
\ddot\gamma^k+\Gamma^k_{ij}\dot\gamma^i\dot\gamma^j=0,
\]
它只是 $\nabla_{\dot\gamma}\dot\gamma=0$ 的坐标写法：\textbf{切向量沿自己平行}。

\subsubsection{（c）内蕴刻画}

球面的高斯曲率 $K=1/R^2$。曲线的\textbf{测地曲率} $k_g$ 度量它在曲面\textbf{内部}看来的弯曲程度。大圆满足 $k_g=0$——它只弯在法方向，不弯在切平面里。

\textbf{测地线 $=$ 测地曲率为零}。这是最本质的定义，因为它不需要嵌入。

\subsection{BR1.3 slerp 公式的推导}

有了测地线，剩下是\textbf{参数化}：要求 $t\in[0,1]$ 与走过的角度成正比（等速）。

设 $\Omega=\arccos\langle p,q\rangle\in(0,\pi)$。在平面 $\mathrm{span}\{p,q\}$ 里令 $\gamma(t)=\alpha(t)p+\beta(t)q$，要求
\[
\langle\gamma(t),p\rangle=\cos(t\Omega),\qquad
\langle\gamma(t),q\rangle=\cos\big((1-t)\Omega\big).
\]
这是关于 $(\alpha,\beta)$ 的 $2\times 2$ 线性方程组，Gram 矩阵为
$\begin{pmatrix}1&\cos\Omega\\ \cos\Omega&1\end{pmatrix}$，解得
\[
\boxed{\ \mathrm{slerp}(p,q;t)=\frac{\sin((1-t)\Omega)}{\sin\Omega}\,p+\frac{\sin(t\Omega)}{\sin\Omega}\,q\ }
\]

\textbf{读法}：分子是“在圆上走弧长”的坐标，分母 $\sin\Omega$ 是归一化因子，保证 $t=0,1$ 精确落在 $p,q$。当 $p\perp q$（$\Omega=\pi/2$）时退化为 $\cos(t\pi/2)p+\sin(t\pi/2)q$。

\subsection{BR1.4 两个奇点，性质完全不同}

\begin{itemize}
  \item $\Omega\to0$：$\sin\Omega\to0$ 是\textbf{坐标奇点}，不是几何奇点。真实几何是“两个旋转几乎相同”，路径是极短的测地线。工程上用分支处理，或改用指数映射写（见 BR1.5）即可消除。
  \item $\Omega\to\pi$：这是\textbf{几何奇点}——对极情形。球面对极两点间有无穷多条等长的大圆，测地线\textbf{不唯一}。无论怎么实现都消灭不了，根源是 $\pi_1(SO(3))=\mathbb{Z}_2$。
\end{itemize}

\textbf{四元数的符号问题}：$q$ 与 $-q$ 是同一旋转。若 $\langle q_0,q_1\rangle<0$，slerp 走的是长弧。所以必须先做
\[
\texttt{if }\langle q_0,q_1\rangle<0:\ q_1\gets -q_1 .
\]
这不是数值技巧，是\textbf{在两个对极代表里挑离得近的那个}——几何上就是选大圆的哪一半。

\subsection{BR1.5 指数映射：slerp 的真名}

在 $p$ 处 $T_pS^n=p^\perp$。定义
\[
\exp_p(v)=\cos\lVert v\rVert\,p+\sin\lVert v\rVert\,\frac{v}{\lVert v\rVert},
\]
这正是 $\ddot\gamma=-\gamma$ 以 $v$ 为初速度的解。于是
\[
\boxed{\ \mathrm{slerp}(q_0,q_1;t)=\exp_{q_0}\big(t\cdot\log_{q_0}(q_1)\big)\ }
\]
其中 $\log_{q_0}(q_1)$ 是“从 $q_0$ 指向 $q_1$、长度为球面距离”的切向量。

\textbf{这一步是本篇关键}：slerp 是\textbf{流形上的指数映射在等速参数下的像}。换到李群语言就是：$S^3=SU(2)$ 上的测地线就是单参数子群 $t\mapsto\exp(tX)$。所以
\[
\text{大圆}\;=\;\text{测地线}\;=\;\text{单参数子群}\;=\;\text{指数映射的像}
\]
这四件事在“紧李群 + 双不变度量”的场合是同一件事。

用四元数乘法写更直白：相对旋转 $r=q_0^{-1}q_1=\cos\frac\theta2+\sin\frac\theta2\hat n$，则
\[
q(t)=q_0\otimes\left(\cos\frac{t\theta}{2}+\sin\frac{t\theta}{2}\hat n\right),
\]
与球面 slerp 同形，因为左乘 $q_0$ 是 $S^3$ 的等距变换，等距把测地线映成测地线。

\subsubsection{等变性：为什么 slerp 是“几何量”}

左乘、右乘都是 $S^3$ 的等距（双不变度量），所以
\[
\mathrm{slerp}(aq_0,aq_1;t)=a\,\mathrm{slerp}(q_0,q_1;t).
\]
\textbf{路径不依赖坐标系}。这是“几何量”与“坐标技巧”的分界：slerp 不是被定义的，是被群作用钉死的。

\subsection{BR1.6 发散：从大圆能牵出什么}

\begin{enumerate}
  \item \textbf{球面距离与余弦相似度}。$d(p,q)=\arccos\langle p,q\rangle$。归一化 embedding 住在球面上，“最像”$=$ 球面距离最近 $=$ 大圆最短。
  \item \textbf{共轭点}。测地线的线性化是 Jacobi 方程
  \[
  \ddot J+R(J,\dot\gamma)\dot\gamma=0 .
  \]
  单位球面上 $R=\mathrm{id}$（切向上），故 $\ddot J+J=0$，解为 $\sin$，在弧长 $\pi R$ 处 $J=0$。\textbf{对极点就是共轭点}。共轭点的意义：到那里测地线不再是唯一极小者（Jacobi 场告诉你可以岔开）——这精确解释了 $\Omega=\pi$ 时 slerp 不唯一。
  \item \textbf{高斯--博内}。三条大圆弧围成的测地三角形满足
  \[
  \alpha+\beta+\gamma=\pi+\int_\triangle K\,dA=\pi+\frac{A}{R^2}.
  \]
  球面上三角形“胖”了，多出的面积正比于曲率。\textbf{曲率的全局指纹}：在球面上做几何，角度的账永远差一块面积。
  \item \textbf{平行移动与 holonomy}。沿大圆把切向量平行搬一圈回来，仍转过 $\int K\,dA$ 的角度。大圆测地曲率为零，是“最不转”的路径，但一圈下来还是转了——球面是弯的。这解释图形学里“绕一圈的旋转不是恒等”（$SO(3)$ 的 holonomy，根源是 $\pi_1(SO(3))=\mathbb{Z}_2$）。
  \item \textbf{常曲率三分}。球面 $K>0$（大圆）、欧氏 $K=0$（直线）、双曲 $K<0$（半圆/直线）。“直线”的样子完全由曲率符号决定，这是黎曼几何第一课。
  \item \textbf{多点插值会出事}。球面不是向量空间，de Casteljau 不能随便搬。三个以上四元数做 spline 要用 squad 或直接在李代数上做。\textbf{原因}：不存在满足所有公理的“球面线性平均”——球面没有线性结构，只有测地中点，而测地中点不满足结合律。多个四元数的“平均”要用 Karcher（Fréchet）均值，迭代求解。
\end{enumerate}
%==============================================================
\section{BR2. 正交基、傅里叶、正态特征函数与拉普拉斯变换}
\label{sec:br2}

\begin{center}
\small
\begin{tabular}{L{3.4cm}L{10.4cm}}
\toprule
\textbf{本篇 $\cong$} & \textbf{一句话} \\
\midrule
傅里叶展开 $\cong$ 正交基上投影 & 选一组让微分算子对角化的基 \\
卷积 $\cong$ 对偶空间里的乘法 & 时域卷积 = 频域逐点乘 \\
正态特征函数 $\cong$ 算子对角化 & $e^{-\sigma^2\omega^2/2}$ 是 $-\partial_x^2$ 的特征值 \\
拉普拉斯变换 $\cong$ 带衰减的内积 & 把 $e^{-st}$ 当基，收敛域决定能不能展开 \\
采样定理 $\cong$ 频域周期延拓 & 混叠 = 延拓后的重叠 \\
\bottomrule
\end{tabular}
\end{center}

\noindent\textbf{落点}：\texttt{H1}--\texttt{H11}（DFT / 卷积定理 / Parseval / 谱收敛）、\texttt{P1}--\texttt{P2}（特征函数与矩）、\texttt{A6}（收敛阶）
%==============================================================

本篇要立一根骨架，然后你会看到四个名词全是它的方言：
\[
\text{一切变换}\;=\;\text{选基}\;+\;\text{对角化},\qquad
f=\sum_k\langle f,e_k\rangle e_k .
\]

\subsection{BR2.1 骨架：为什么偏偏选复指数}

内积空间里若 $\{e_k\}$ 正交归一（$\langle e_j,e_k\rangle=\delta_{jk}$），则
\[
f=\sum_k\langle f,e_k\rangle e_k .
\]
\textbf{系数就是内积，没有别的。} 对照泰勒展开：$\{1,x,x^2,\dots\}$ 不正交，系数只能靠求导，很别扭。

\textbf{为什么选复指数}：
\[
\frac{d}{dx}e^{inx}=in\,e^{inx},
\]
即\textbf{在傅里叶基下，微分算子是对角的}（$d/dx\leftrightarrow\mathrm{diag}(in)$）。这就是傅里叶变换能把微分方程变成代数方程的真正原因——不是巧合，是选对了基。

\subsubsection{Sturm--Liouville 框架}

任何二阶自伴算子
\[
L=-\frac{d}{dx}\left(p\frac{d}{dx}\right)+q
\]
配上合适边界条件，都有完备正交本征函数系，$L\phi_n=\lambda_n w\phi_n$。\textbf{选哪个基，取决于解哪个算子、在哪个区域}：

\begin{center}
\begin{tabular}{L{4.2cm}L{3.2cm}L{4.5cm}}
\toprule
\textbf{算子 / 区域} & \textbf{本征函数} & \textbf{出现场合} \\
\midrule
$d^2/dx^2$，周期 & $e^{inx}$ & 傅里叶级数 \\
球坐标 $\theta$ 方向 & 勒让德 $P_l$ & 球谐、氢原子 \\
柱坐标 & 贝塞尔 $J_\nu$ & 圆膜振动、热传导 \\
$d^2/dx^2-x^2$（谐振子） & Hermite $H_ne^{-x^2/2}$ & \textbf{傅里叶变换的本征函数} \\
区间 $(0,\infty)$ & 拉盖尔 & 氢原子径向 \\
\bottomrule
\end{tabular}
\end{center}

\textbf{倒数第二行是整篇的伏笔}：Hermite 函数满足
\[
\mathcal{F}[\psi_n]=(-i)^n\psi_n,\qquad n=0:\ \ \mathcal{F}\left[e^{-x^2/2}\right]=e^{-\omega^2/2}.
\]
\textbf{高斯是傅里叶变换的不动点（本征值 1）}。后面所有“正态为什么特殊”的问题，答案都是它。

\subsection{BR2.2 傅里叶级数到傅里叶变换：只是基变连续}

周期 $T$ 的函数：
\[
f(x)=\sum_n c_ne^{i\omega_n x},\quad
c_n=\frac1T\int_{-T/2}^{T/2}f\,e^{-i\omega_n x}dx,\quad
\omega_n=\frac{2\pi n}{T}.
\]
令 $T\to\infty$：频率间隔 $\Delta\omega=2\pi/T\to 0$，求和变积分：
\[
f(x)=\frac{1}{2\pi}\int\hat f(\omega)e^{i\omega x}d\omega,\qquad
\hat f(\omega)=\int f(x)e^{-i\omega x}dx .
\]
严谨化需要 Plancherel 定理（$L^2$ 等距同构）与谱定理。

\subsubsection{Pontryagin 对偶：级数与变换不是两种东西}

傅里叶变换是\textbf{平移群表示的“对角化”}，特征标是 $e^{i\omega x}$。而
\[
\widehat{\mathbb{R}}\cong\mathbb{R},\qquad
\widehat{\mathbb{Z}}\cong S^1,\qquad
\widehat{S^1}\cong\mathbb{Z}.
\]
\textbf{周期函数就是 $S^1$ 上的函数，它的对偶是 $\mathbb{Z}$，所以傅里叶级数的系数是“$\mathbb{Z}$ 上的函数”（离散谱）}。级数与变换是同一对偶框架在两个群上的样子，这就是“离散谱与连续谱”的终极解释。

\subsection{BR2.3 特征函数：概率论里的傅里叶变换}

\[
\varphi_X(t)=\mathbb{E}[e^{itX}]=\int e^{itx}dF(x).
\]
\textbf{傅里叶变换的所有定理，立刻变成概率论定理}：

\begin{center}
\begin{tabular}{L{4.5cm}L{6.5cm}}
\toprule
\textbf{傅里叶} & \textbf{概率} \\
\midrule
变换是单射 & $\varphi_X$ 唯一决定分布 \\
卷积 $\leftrightarrow$ 乘积 & $X\perp Y$ 时 $\varphi_{X+Y}=\varphi_X\varphi_Y$ \\
求导 $\leftrightarrow$ 乘 $x$ & $\varphi^{(k)}(0)=i^k\mathbb{E}[X^k]$ \\
反转公式 & $f(x)=\frac{1}{2\pi}\int e^{-itx}\varphi(t)dt$ \\
高斯自对偶 & $N(0,1)$ 的 $\varphi$ 与密度同形 \\
\bottomrule
\end{tabular}
\end{center}

\subsection{BR2.4 正态特征函数的严格计算}

$X\sim N(0,1)$：
\[
\varphi(t)=\frac{1}{\sqrt{2\pi}}\int e^{-x^2/2+itx}dx .
\]
配方：$-x^2/2+itx=-\frac12(x-it)^2-\frac{t^2}{2}$，于是
\[
\varphi(t)=e^{-t^2/2}\cdot\frac{1}{\sqrt{2\pi}}
\int_{\mathbb{R}}e^{-\frac12(x-it)^2}dx .
\]

\textbf{关键不是配方，是“为什么可以把路径从实轴移到 $\mathrm{Im}\,z=-t$”}：

取矩形围道（上下边 $\mathrm{Im}=0$ 与 $\mathrm{Im}=-t$，两侧竖边 $\pm R$）。$e^{-z^2/2}$ 在 $\mathbb{C}$ 上\textbf{整}，矩形内积分为零。两侧竖边上
$|e^{-z^2/2}|=e^{-(R^2-y^2)/2}\to0$（$|y|\le|t|$ 一致），故\textbf{上下两边积分相等}：
\[
\int_{\mathbb{R}}e^{-\frac12(x-it)^2}dx=\int_{\mathbb{R}}e^{-x^2/2}dx=\sqrt{2\pi}.
\]
所以 $\varphi(t)=e^{-t^2/2}$。一般 $N(\mu,\sigma^2)$：
\[
\boxed{\ \varphi(t)=\exp\left(i\mu t-\frac{\sigma^2t^2}{2}\right)\ }
\]
它和密度函数\textbf{形状完全相同}——这是 BR2.1 那条 Hermite 本征性质的直接后果。

\subsection{BR2.5 中心极限定理：用 Taylor 展开看“为什么是高斯”}

$X_i$ iid、均值 $0$、方差 $1$，$S_n=\frac{1}{\sqrt n}\sum X_i$：
\[
\varphi_{S_n}(t)=\left[\varphi\left(\frac{t}{\sqrt n}\right)\right]^n .
\]
原点附近 $\varphi(t)=1-\frac{t^2}{2}+o(t^2)$，故
\[
\varphi_{S_n}(t)=\left(1-\frac{t^2}{2n}+o\left(\frac1n\right)\right)^n\longrightarrow e^{-t^2/2}.
\]
再由 Lévy 连续性定理（特征函数逐点收敛 $\Rightarrow$ 分布弱收敛）得 $S_n\Rightarrow N(0,1)$。

\textbf{本质}：所有分布的特征函数在原点附近都长得像 $1-t^2/2$，而它就是高斯特征函数对数的一阶近似。取 $n$ 次幂，高阶项被 $1/n$ 压死，只剩高斯。

\subsubsection{隐藏条件：方差必须有限}

若方差无限（如 Cauchy），$\varphi$ 在原点附近不是 $1-t^2/2$ 而是 $1-c|t|^\alpha$（$\alpha<2$），收敛到\textbf{稳定分布}：
\[
\varphi(t)=\exp\left(i\mu t-c|t|^\alpha\left(1-i\beta\,\mathrm{sgn}(t)\,\Phi(t,\alpha)\right)\right),
\]
$\alpha=2$ 是高斯，$\alpha=1$ 是 Cauchy。\textbf{“高斯特殊不特殊”的答案是：它是所有稳定分布里唯一二阶矩有限的（$\alpha=2$），因此也是 CLT 唯一的归宿}。这从背面把高斯钉死。

\subsection{BR2.6 拉普拉斯变换：给傅里叶加一扇窗}

\[
\mathcal{L}[f](s)=\int_0^\infty f(t)e^{-st}dt,\qquad s=\sigma+i\omega .
\]
精确关系：
\[
\mathcal{L}[f](\sigma+i\omega)=\mathcal{F}\left[f(t)e^{-\sigma t}\mathbf{1}_{t>0}\right](\omega).
\]
\textbf{拉普拉斯 $=$ 单边 $+$ 乘指数衰减窗的傅里叶变换}。$e^{-\sigma t}$ 把指数增长的 $f$ 压成可积；$\sigma$ 的自由度就是收敛域（ROC）。

\subsubsection{为什么工程上更爱拉普拉斯}

\begin{enumerate}
  \item \textbf{初值自动进门}：$\mathcal{L}[f']=sF(s)-f(0)$。求导变乘 $s$，初值变成加项，ODE 初值问题直接变代数方程——传递函数 $H(s)$ 的来源。
  \item \textbf{因果性}：$t>0$ 才有，卷积是 $[0,\infty)$ 上的卷积。
  \item \textbf{极点即模态}：极点 $\pm i\omega_0$ 等幅振荡；实部 $<0$ 衰减；实部 $>0$ 发散。\textbf{稳定性判据变成“极点在哪半个平面”}。
  \item \textbf{三件套}：$s=i\omega$ 是稳态（傅里叶），$z=e^{sT}$ 是离散（$z$ 变换），$s$ 平面是连续频域。
\end{enumerate}

\subsubsection{与矩生成函数的关系}

矩生成函数 $M(t)=\mathbb{E}[e^{tX}]$，特征函数 $\varphi(t)=\mathbb{E}[e^{itX}]$，故
\[
\varphi(t)=M(it).
\]
\textbf{同一个函数在实轴与虚轴上的两副面孔}。正态的
\[
M(t)=\exp\left(\mu t+\frac{\sigma^2t^2}{2}\right),\qquad
\varphi(t)=\exp\left(i\mu t-\frac{\sigma^2t^2}{2}\right)
\]
正好是 $t\to it$ 的解析延拓。\textbf{正态是唯一让 MGF 与 CF 都闭式、且互为解析延拓的那一档}（因为它的对数在 $t$ 上是二次的，二次函数在旋转下闭）。

\subsection{BR2.7 不确定性原理与“最集中”}

正交基的完备性 + Plancherel 等距性推出：$f$ 与 $\hat f$ 不能同时窄。
\[
\sigma_x\,\sigma_\omega\ge\frac12,
\]
等号\textbf{恰好}在高斯时成立。检验：$f=e^{-x^2/4\sigma^2}$，$\hat f$ 宽度 $1/\sigma$，乘积为常数。

\textbf{三者锁在一起}：正交基完备 $\Rightarrow$ 不确定性 $\Rightarrow$ 高斯自对偶。

\subsection{BR2.8 三处补充}

\subsubsection{（一）热核就是高斯}

热方程 $u_t=\Delta u$，傅里叶变换后 $\hat u_t=-|\xi|^2\hat u$，故
$\hat u(\xi,t)=e^{-t|\xi|^2}\hat u_0(\xi)$。逆变换（高斯自对偶 + 缩放）：
\[
u(x,t)=(G_t*u_0)(x),\qquad G_t(x)=(4\pi t)^{-n/2}e^{-|x|^2/4t}.
\]
\textbf{热核就是高斯}。“抹平初始条件”的高斯刷子，在频域里是乘 $e^{-t|\xi|^2}$——一个对角算子。这是“正交基把算子对角化”最实用的一次。

\subsubsection{（二）小波：局部化的正交基}

傅里叶基在空间上无穷宽、在频率上无穷窄（一条线）；小波基相反，是\textbf{时频平面铺砖}。两者同为 $L^2$ 正交基，差别只是砖怎么切。\textbf{选基的第三条依据}：除对称性与算子，还有“你关心的东西是局部还是全局”。

\subsubsection{（三）Mercer 核与 KL 展开}

正定核
\[
K(x,y)=\sum_k\lambda_k\phi_k(x)\overline{\phi_k(y)}
\]
是协方差算子的谱分解。高斯核对应的 $\phi_k$ 是 Hermite 函数（Mehler 核），所以\textbf{高斯核也是“高斯”}——这是高斯过程处处见闭式解的原因。PCA 是有限维版本的同一件事。
%==============================================================
\section{BR3. ODE、PDE、张量、李群：一张网}
\label{sec:br3}

\begin{center}
\small
\begin{tabular}{L{3.4cm}L{10.4cm}}
\toprule
\textbf{本篇 $\cong$} & \textbf{一句话} \\
\midrule
ODE $\cong$ 向量场的积分 & 解就是流的轨道 \\
PDE $\cong$ 无穷维向量场的流 & 空间离散化之后回到 ODE \\
张量 $\cong$ 多线性映射 & 换基就是坐标变换，指标只是记账 \\
李群 $\cong$ 带群结构的流形 & 群运算本身光滑 \\
李代数 $\cong$ 单位元处的切空间 & 在恒等元处局部线性化 \\
外微分 $\cong$ 反对称化的导数 & $d^2=0$ 就是「边界的边界为空」 \\
\bottomrule
\end{tabular}
\end{center}

\noindent\textbf{落点}：\texttt{IF-C1}/\texttt{IF-C2}（ode / forms）、\texttt{W1}、\texttt{T1}--\texttt{T7}（张量）、\texttt{M1}--\texttt{M7}（微分形式）、\texttt{E-ODE-3}、\texttt{D6}--\texttt{D8}
%==============================================================

\noindent\small\textbf{这里默认你已经会}：一阶/二阶 ODE 的解法与初值问题；多元微积分
（偏导、链式法则、Jacobian、Hessian、重积分）；矩阵指数 $e^{At}$ 的级数定义；
基变换时分量的变化规则；群的四条定义。\\
\textbf{不要求会}：Christoffel 符号的计算、黎曼曲率的推导、李代数表示论、BCH 的证明 ——
全篇只给结论、不给推导。\textbf{这是四篇里门槛最高的一篇}：\S BR3.3 的曲率部分与
\S BR3.4 后半最容易卡，跳过它们不影响 \S BR3.6 的读法。\normalsize
\smallskip

本篇不追求定理完整，追求\textbf{四者之间的接线}。先给主图：

\begin{center}
\begin{tabular}{L{2.2cm}L{9.5cm}}
\toprule
\textbf{层} & \textbf{内容} \\
\midrule
对称 & 连续对称（李群 $G$） \\
$\downarrow$ 无穷小 & 李代数 $\mathfrak{g}=T_eG$，带括号 $[\cdot,\cdot]$ \\
$\downarrow$ 指数映射 & 流 $\Phi_t$（ODE 的解，单参数群） \\
$\downarrow$ 几何量 & 张量场（$g$、$R$、应力） \\
$\downarrow$ 局部分程 & PDE（Laplace--Beltrami、NS、Einstein） \\
\bottomrule
\end{tabular}
\end{center}

\subsection{BR3.1 ODE：流的语言}

核心方程 $\dot x=f(t,x)$，$x$ 在 $\mathbb{R}^n$ 或流形 $M$ 上。

\subsubsection{存在唯一与流}

\textbf{Picard--Lindelöf}：$f$ 局部 Lipschitz $\Rightarrow$ 局部唯一解。证明是不动点迭代
\[
x_{k+1}(t)=x_0+\int_0^t f(s,x_k(s))ds,
\]
几何意义就是“沿向量场走无穷小步”——它与指数映射是同一种东西的两种写法。

\textbf{流} $\Phi_t:x_0\mapsto x(t)$ 满足
\[
\Phi_t\circ\Phi_s=\Phi_{t+s},\qquad \Phi_0=\mathrm{id}.
\]
\textbf{这是单参数群}。所以 ODE 天然生成群结构——李群不是外加的，是 ODE 的必然产物。

\subsubsection{线性系统与矩阵指数}

$\dot x=Ax$ 的解 $x(t)=e^{At}x_0$，
\[
e^{At}=\sum_{k\ge0}\frac{(At)^k}{k!}.
\]
Jordan 型 $\Rightarrow$ 稳定性看 $\mathrm{Re}\,\lambda$。\textbf{$e^{At}$ 就是 $GL(n)$ 的指数映射}——ODE 与李群的第一个硬连接。

\textbf{平衡点与线性化}：$\dot x=Ax+O(|x|^2)$。双曲平衡点附近 Hartman--Grobman：非线性流与线性流拓扑共轭。稳定/不稳定方向的维数由 $\mathrm{Re}\,\lambda$ 的符号决定（稳定流形定理）。

\textbf{Lyapunov}：$V>0$，$\dot V=\nabla V\cdot f\le0$ $\Rightarrow$ 稳定。$V$ 是能量，$\dot V$ 是能量沿流的变化率。

\subsubsection{Hamilton 系统：辛几何的入口}

\[
\dot q=\frac{\partial H}{\partial p},\qquad
\dot p=-\frac{\partial H}{\partial q}
\quad\Longleftrightarrow\quad
\dot z=J\nabla H,
\]
其中 $J=\begin{pmatrix}0&I\\-I&0\end{pmatrix}$，$J^2=-I$，$J^T=-J$。

\textbf{这里出现了辛结构}：$J$ 是反对称、满足 $J^2=-I$ 的“复结构”。Hamilton 流保辛形式 $\omega=\sum dq_i\wedge dp_i$（Liouville 定理）。\textbf{辛几何的入口在 ODE 里，不在几何书里。}

\subsubsection{边值问题给出谱}

Sturm--Liouville
\[
-(py')'+qy=\lambda wy
\]
本征值离散、本征函数正交完备。\textbf{这就是 BR2.1“选基”的数学来源}：基不是天上掉的，是边值问题的本征函数。

\subsection{BR3.2 PDE：把算子放到函数空间上}

方程 $F(x,u,Du,D^2u)=0$。注意 $Du$ 是 Jacobian（一阶张量），$D^2u$ 是 Hessian（二阶对称张量）——\textbf{PDE 从一开始就是张量方程}。

\subsubsection{二阶线性分类}

$au_{xx}+2bu_{xy}+cu_{yy}+\cdots=0$，判别式 $\Delta=b^2-ac$：

\begin{center}
\begin{tabular}{L{1.6cm}L{2.2cm}L{3.2cm}L{2.6cm}L{3.4cm}}
\toprule
$\Delta$ & 类型 & 代表 & 定解条件 & 解的性质 \\
\midrule
$<0$ & 椭圆 & $\Delta u=0$ & 边值 & 光滑、极大值原理 \\
$=0$ & 抛物 & $u_t=\Delta u$ & 初值+边值 & 瞬时抹平、半群 \\
$>0$ & 双曲 & $u_{tt}=\Delta u$ & 初值 & 特征线、有限传播速度 \\
\bottomrule
\end{tabular}
\end{center}

\textbf{分类的几何意义 $=$ 主符号的二次型符号} $a\xi^2+2b\xi\eta+c\eta^2$。换成张量语言，就是\textbf{二阶对称张量 $a_{ij}$ 的惯性指数}。“椭圆/抛物/双曲”不是三个孤立名字，是同一二次型在三种符号下的三种命运。

\subsubsection{三大方程}

\begin{enumerate}
  \item \textbf{Laplace} $\Delta u=0$：均值性质、极大值原理、Green 函数（基本解 $-\frac{1}{2\pi}\log r$ 或 $\frac{1}{4\pi r}$）。
  \item \textbf{热} $u_t=\Delta u$：$u(t)=e^{t\Delta}u_0$，\textbf{半群}；傅里叶解得核为高斯（BR2.8）。
  \item \textbf{波} $u_{tt}=\Delta u$：$\hat u=e^{\pm it|\xi|}\hat u_0$；一维 d'Alembert $u=f(x-t)+g(x+t)$，\textbf{特征线}。
\end{enumerate}

\subsubsection{变分原理：把 PDE 变回几何}

\begin{itemize}
  \item Laplace $=$ Dirichlet 能量 $E(u)=\frac12\int|\nabla u|^2$ 的临界点；
  \item 波动 $=$ 最小作用量；
  \item \textbf{PDE 是泛函的 Euler--Lagrange 方程}。这把 PDE 从“偏微分方程组”变成“几何的驻值条件”。
\end{itemize}

\subsubsection{弱形式：内积的语言}

\[
\int_\Omega u\,\Delta v=-\int_\Omega\nabla u\cdot\nabla v+\int_{\partial\Omega}u\,\partial_n v .
\]
乘测试函数再分部积分，二阶导降为一阶——弱解与有限元的起点。\textbf{它和 BR2.1 的“内积展开”是同一件事}：用内积把方程投影到测试函数上。

\textbf{特征线}：$\partial_tu+a\partial_xu=0$ 沿 $x-at=$ 常数不变。\textbf{一阶 PDE $=$ 一族 ODE}。非线性时特征线相交 $\Rightarrow$ 激波（有限时间形成奇性）。这是“PDE 的困难来自非线性”的第一个实例。

\subsection{BR3.3 张量：坐标无关的多线性代数}

\subsubsection{三种视角}

\begin{enumerate}
  \item \textbf{分量}：带上下标的数，按变换律变。逆变 $T'^i=\frac{\partial x'^i}{\partial x^k}T^k$；协变 $T'_j=\frac{\partial x^k}{\partial x'^j}T_k$。
  \item \textbf{多重线性映射}：$T:V^*\times\cdots\times V^*\times V\times\cdots\to\mathbb{R}$，对每个槽线性。
  \item \textbf{张量积}：$T\in V\otimes\cdots\otimes V^*\otimes\cdots$，基是 $e_i\otimes\cdots\otimes e^k$。
\end{enumerate}

\subsubsection{为什么需要它}

因为物理/PDE 定律必须\textbf{坐标无关}。在球坐标写 NS，与在直角坐标写 NS，必须是同一个物理。张量是“保证换坐标后方程形式不变”的语法。\textbf{张量不是数学家的洁癖，是“换坐标系不会把答案换掉”的保险。}

\subsubsection{关键操作}

\begin{itemize}
  \item \textbf{缩并}：$T^i{}_i$（一上一下配对求和）。这是张量代数里唯一自然“降秩”的操作。
  \item \textbf{对称/反对称分解}：任意二阶张量 $=$ 对称 $+$ 反对称；3 维中反对称部分对应轴向量 $v_i=\epsilon_{ijk}A^{jk}$。
  \item \textbf{度量}：升指标 $T^i=g^{ij}T_j$、降指标 $T_j=g_{ij}T^i$。度量定义长度、角度、体积元 $\sqrt{|g|}$。
  \item \textbf{协变导数}：
  \[
  \nabla_kT^i{}_j=\partial_kT^i{}_j+\Gamma^i_{kl}T^l{}_j-\Gamma^l_{kj}T^i{}_l ,
  \]
  $\Gamma$ 是 Christoffel 符号，由 $g$ 唯一确定（无挠 $+$ 度量相容）。
  \item \textbf{曲率 $=$ 协变导数的不对易性}：
  \[
  [\nabla_k,\nabla_l]T^i=R^i{}_{jkl}T^j .
  \]
\end{itemize}

\subsubsection{黎曼曲率的对称性}

$R_{ijkl}=g_{im}R^m{}_{jkl}$ 满足
\[
R_{ijkl}=-R_{jikl}=-R_{ijlk}=R_{klij},\qquad
R_{ijkl}+R_{iklj}+R_{iljk}=0 .
\]
独立分量数 $\dfrac{n^2(n^2-1)}{12}$：$n=2$ 时 1 个（就是高斯曲率），$n=3$ 时 6 个，$n=4$ 时 20 个。

里奇 $R_{ij}=R^k{}_{ikj}$，标量曲率 $R=g^{ij}R_{ij}$，爱因斯坦张量 $G_{ij}=R_{ij}-\frac12Rg_{ij}$（$\nabla^iG_{ij}=0$ 是 Bianchi 恒等式的缩写）。

\subsubsection{高斯--博内：接回 BR1}

\[
\int_M K\,dA=2\pi\chi(M)
\]
（2 维闭曲面）。球面 $\chi=2$，总曲率 $4\pi$。\textbf{球面三角形内角和 $>\pi$ 与 slerp 走大圆，是同一枚硬币}：在弯曲空间做几何，角度和面积的账永远差一个曲率积分。

\subsubsection{张量进 PDE}

\begin{itemize}
  \item \textbf{Laplace--Beltrami}：
  \[
  \Delta_gu=\frac{1}{\sqrt{|g|}}\partial_i\left(\sqrt{|g|}g^{ij}\partial_ju\right).
  \]
  球坐标下即 $\frac{1}{r^2}\partial_r(r^2\partial_r)+\frac{1}{r^2}\Delta_{S^2}$，$S^2$ 部分就是球谐算子。
  \item \textbf{NS 的应力形式}：$\sigma_{ij}=-p\delta_{ij}+2\mu e_{ij}$，$e_{ij}=\frac12(\partial_iu_j+\partial_ju_i)$，平衡方程 $\nabla_j\sigma^{ij}=f^i$。\textbf{NS 的几何写法就是“应力张量的散度等于外力”。}
  \item \textbf{Einstein 方程} $G_{ij}=8\pi T_{ij}$：10 个二阶非线性 PDE，未知量是 $g_{ij}$。与 NS 同一套路——一个张量（几何）等于另一个张量（物质）。
\end{itemize}

\textbf{微分形式 $=$ 反对称张量}。外积 $\wedge$、外微分 $d$（$d^2=0$）。Hodge 分解
\[
L^2(\Lambda^k)=d(\Lambda^{k-1})\oplus\delta(\Lambda^{k+1})\oplus\mathcal{H}^k ,
\]
3 维上即 Helmholtz 分解：\textbf{无旋 $+$ 无散 $+$ 调和}。\textbf{这就是不可压条件 $\nabla\cdot u=0$ 的几何本质}——$u$ 落在 $\delta(\Lambda^2)$ 那一块，压力是投影到无散子空间的 Lagrange 乘子。

\subsection{BR3.4 李群：连续对称的语言}

\subsubsection{经典李群}

\begin{center}
\begin{tabular}{L{2.8cm}L{1.6cm}L{6cm}}
\toprule
\textbf{群} & \textbf{维数} & \textbf{说明} \\
\midrule
$GL(n)$ & $n^2$ & 可逆矩阵 \\
$SL(n)$ & $n^2-1$ & $\det=1$ \\
$O(n)$ & $n(n-1)/2$ & 保长度 \\
$SO(n)$ & $n(n-1)/2$ & $\det=1$，$SO(3)$ 是 3 维旋转 \\
$SU(n)$ & $n^2-1$ & $\det=1$ 的酉群 \\
$S^3\cong SU(2)$ & 3 & \textbf{单位四元数} \\
$SE(3)$ & 6 & 刚体运动（旋转 $+$ 平移） \\
\bottomrule
\end{tabular}
\end{center}

\textbf{关键拓扑}：$SO(3)\cong SU(2)/\{\pm I\}$ 是\textbf{双重覆盖}，$\pi_1(SO(3))=\mathbb{Z}_2$。$SU(2)$ 中 $\theta:0\to2\pi$ 走到 $-I$，\textbf{不是恒等}；走两圈才回 $I$。这就是 BR1.4“$q$ 与 $-q$ 是同一旋转”的拓扑根源。

\subsubsection{李代数与指数映射}

$\mathfrak{g}=T_eG$，带括号 $[X,Y]$；矩阵群即 $[X,Y]=XY-YX$。\textbf{李代数 $=$ 群在单位元的线性化}，它丢掉全局拓扑（$SU(2)$ 与 $SO(3)$ 的李代数都是 $\mathfrak{so}(3)$，但群不同）——这是“$SU(2)$ 能表示半整数自旋而 $SO(3)$ 不能”的准确说法。

指数映射 $\exp:\mathfrak{g}\to G$ 是过 $e$、初速 $X$ 的积分曲线在 $t=1$ 的值。

\begin{itemize}
  \item $\mathfrak{su}(2)\to SU(2)$：$\exp(v)=\cos\lVert v\rVert+\frac{\sin\lVert v\rVert}{\lVert v\rVert}v$ ——\textbf{这就是 BR1 的 slerp}。
  \item $\mathfrak{so}(3)\to SO(3)$：Rodrigues $R=I+\sin\theta\,K+(1-\cos\theta)K^2$，$K$ 是单位轴的叉积矩阵。
\end{itemize}

\textbf{测地线 $=$ 单参数子群}：双不变度量下，李群上的测地线恰是 $t\mapsto\exp(tX)$。\textbf{BR1 与 BR3 在这里握手}。

\subsubsection{BCH 公式：为什么旋转不能相加}

\[
\log(e^Xe^Y)=X+Y+\frac12[X,Y]+\frac{1}{12}[X,[X,Y]]-\frac{1}{12}[Y,[X,Y]]+\cdots
\]
二阶项是交换子。\textbf{slerp 分两步与一步的差就是 $\frac12[X,Y]$}。非交换性不是抽象代数的事，是两次小旋转叠加时那个不该被忽略的误差。

\subsubsection{伴随与表示}

$\mathrm{Ad}_gX=gXg^{-1}$，$\mathrm{ad}_X=[X,\cdot]$。$SU(2)$ 不可约表示由 $j=0,\frac12,1,\dots$ 标号，维数 $2j+1$；$SO(3)$ 只有整数 $j$。

\textbf{闭环到 BR2}：$SO(3)$ 作用在 $\mathbb{R}^3$ 上的张量按不可约表示分解
\[
\mathbb{R}^3\otimes\mathbb{R}^3=\underbrace{\mathbb{R}}_{\text{迹}}\oplus\underbrace{\Lambda^2}_{\text{反对称}}\oplus\underbrace{S_0}_{\text{对称无迹}},
\]
对应 $j=0,1,2$。\textbf{球谐 $Y_{lm}$ 就是这个分解在 $L^2(S^2)$ 上的展开基}。所以球坐标下解 Laplace 自动得到球谐，不是巧合，是群论——“选基”最终是群论的事。

\textbf{Noether}：作用量对连续对称不变 $\Rightarrow$ 守恒流。时间平移 $\to$ 能量，空间平移 $\to$ 动量，$U(1)$ 相位 $\to$ 电荷。\textbf{这是李群进物理的唯一入口，也是“物理定律总写成张量形式”的答案}。

\textbf{刚体 $=$ $SO(3)$ 上的测地线}：Euler 方程 $\dot L=L\times\omega$ 就是 $SO(3)$ 带惯性度量的测地线方程。\textbf{角动量守恒在这里是“测地线 $=$ 单参数子群”的特例}——ODE、张量、李群、几何全焊在一处。

\subsection{BR3.5 交叉点}

\begin{enumerate}
  \item \textbf{ODE $+$ 李群}：$e^{At}$ 是 $GL(n)$ 的单参数子群；Hamilton 流是辛群的子群；刚体方程是 $SO(3)$ 的测地线。
  \item \textbf{ODE $+$ 张量}：平行移动 $\nabla_{\dot\gamma}T=0$ 是一阶线性 ODE；测地线方程 $\ddot\gamma^k+\Gamma^k_{ij}\dot\gamma^i\dot\gamma^j=0$ 就是“$\nabla_{\dot\gamma}\dot\gamma=0$”的坐标展开。
  \item \textbf{张量 $+$ PDE}：Laplace--Beltrami、应力张量散度、Einstein 方程。
  \item \textbf{PDE $+$ 李群}：对称性约化（球对称把 3 维压成 1 维 ODE）；Noether 守恒；Burgers $u_t+uu_x=\nu u_{xx}$ 经 Cole--Hopf $u=-2\nu(\log\theta)_x$ 变成热方程——\textbf{本质是“非线性变换把方程拉直”}，与群作用下的等价变换同类。
  \item \textbf{李群 $+$ 张量 $+$ BR2}：$SO(3)$ 不可约表示 $\to$ 球谐 $\to$ 球坐标下的正交基。
\end{enumerate}

\subsection{BR3.6 用 NS 收尾}

\[
\partial_tu+(u\cdot\nabla)u=-\nabla p+\nu\Delta u,\qquad \nabla\cdot u=0 .
\]
\begin{itemize}
  \item $\nabla\cdot u=0$ 用 Hodge 分解理解：投影到无散子空间。令 $P$ 为 Leray 投影，方程变为
  \[
  \partial_tu+\nu Au+B(u,u)=0,\qquad A=-P\Delta,\ \ B(u,u)=P(u\cdot\nabla u).
  \]
  \item $A$ 在 $\mathcal{H}=\{u\in L^2:\nabla\cdot u=0,\ u\cdot n|_{\partial\Omega}=0\}$ 上\textbf{自伴正定}，故有正交基（Stokes 特征函数）。\textbf{“解 NS 用什么基”的答案就是 $\mathcal{H}$ 的谱分解}——BR2.1 的框架第四次出现。
  \item $B$ 双线性且反对称（$\langle B(u,u),u\rangle=0$，能量守恒）。这是 2 维能证全局存在（能量估计封死）、3 维卡在涡量拉伸 $\omega\cdot\nabla u$ 的原因——3 维没有 2 维那样的守恒律。
  \item 定常 NS 的几何形式就是 $\nabla_j\sigma^{ij}=f^i$（应力张量）。
\end{itemize}

\subsection{BR3.7 入门够用的五问}

\begin{enumerate}
  \item 为什么球面上最直的路是大圆？（$\ddot\gamma=-\gamma$）
  \item 为什么正态特征函数还是高斯形状？（Hermite $+$ 傅里叶自对偶）
  \item 为什么拉普拉斯能处理初值而傅里叶不能？（单边 $+$ 指数窗）
  \item 为什么不可压 NS 要先做 Hodge 分解？（约束的 Lagrange 乘子是压力）
  \item 为什么 $q$ 与 $-q$ 是同一旋转，插值又不能随便选？（双重覆盖 $+$ 对极是共轭点）
\end{enumerate}
%==============================================================
\section{BR4. 概率论的双向地图}
\label{sec:br4}
%==============================================================

本篇的组织方式是“层”：每层给\ \textbf{数学对象 $\to$ 概率对应 $\to$ 反馈回数学的产物}。
先说一句结论，它决定了整篇的读法：

\begin{quote}
\textbf{概率论的原创概念只有三个：独立性、条件化、鞅。}
测度论一个都没有。所以概率论不是“测度论的应用”，
它是“在测度论上装了一根时间轴和一根信息轴”之后长出来的新学科。
\end{quote}

\subsection{BR4.1 五层总表}

\begin{center}
\begin{tabular}{L{1.5cm}L{4.2cm}L{4.2cm}L{4.3cm}}
\toprule
\textbf{层} & \textbf{数学 $\to$ 概率} & \textbf{概率 $\to$ 数学} & \textbf{最易错处} \\
\midrule
0 测度 & 测度 $\to P$，积分 $\to E$ & RN 导数、五种收敛、胎紧 & 独立性被降级成“乘积测度”一格 \\
1 泛函 & $L^p$、投影、半群 & 条件期望、鞅空间、$H^1$-$BMO$ & “弱收敛”实为 weak-$*$ \\
2 调和 & 傅里叶 $\to$ 特征函数 & 乘子、L\'evy-Khintchine & Bochner 是“正定函数”那个定理 \\
3 几何 & $g^{-1}\to\Sigma$、$\Delta_g$ & $\Gamma_2$ 判据、$W_2$、集中 & 曲率的桥是 $\Gamma,\Gamma_2$，不是 Ric \\
4 PDE & 椭圆算子 $\to$ 生成元 & Feynman--Kac、大偏差 & 枢纽是 Legendre 变换 \\
5 代数 & 群/图/素数 $\to$ 随机模型 & 随机矩阵、概率方法 & 自由独立是第三种独立性 \\
\bottomrule
\end{tabular}
\end{center}

\subsection{BR4.2 第 0 层：测度 $\to$ 概率空间}

\subsubsection{对应与精确条件}

\begin{center}
\begin{tabular}{L{3.6cm}L{3.4cm}L{6.5cm}}
\toprule
\textbf{数学} & \textbf{概率} & \textbf{要小心} \\
\midrule
$(S,\mathcal{A})$ & $(\Omega,\mathcal{F})$ & $\mathcal{F}$ 是 $\sigma$-代数 \\
测度 $\mu$ & $P$，$P(\Omega)=1$ & \textbf{总质量 1 是唯一新东西} \\
可测函数 $f$ & 随机变量 $X$ & $X^{-1}(B)\in\mathcal{F}$ \\
$\int f\,d\mu$ & $\mathbb{E}[X]$ & 同一个积分 \\
零测集 & a.s. & \textbf{只对可数多个条件同时成立} \\
$d\mu/d\nu$ & 密度 $p=dP/d\lambda$ & \textbf{需 $P\ll\lambda$ 且 $\sigma$-有限} \\
$\mu\otimes\nu$ & \textbf{独立性} & 概率论唯一的原创概念 \\
条件测度 & 条件分布 & 需正则性（Polish 空间保证） \\
\bottomrule
\end{tabular}
\end{center}

\subsubsection{三个要点}

\textbf{（a）“总质量 1”不是小条件。} 它让 $\mathcal{M}_1(\Omega)$ 成为\textbf{凸集}：凸组合 $\alpha P+(1-\alpha)Q$ 仍是概率测度，于是“混合分布”有几何意义，Fisher 度量、信息几何全部从这个凸结构长出来。\textbf{“1”这个数字的后果远大于“归一化”。}

\textbf{（b）a.s. 的可数性。} $P(A_n)=1$ 对每个 $n$ 成立 $\Rightarrow P(\bigcap_n A_n)=1$（零测集的可数并仍零测）。\textbf{但不可数交不一定。}

后果：像“对几乎所有 $\omega$、对所有 $t\ge0$，$B_t$ 连续”这句话，\textbf{不能}从“对每个 $t$，$B_t$ 连续（a.s.）”直接推出——因为 $t$ 不可数。Brownian motion 的连续性必须先选一个\textbf{可分离版本}（Kolmogorov 连续性定理的用途）。

一句话：\textbf{a.s. 不是“处处”，它的例外只能在可数多个条件下被同时控制。} 这是随机过程的第一道门槛。

\textbf{（c）独立性 $=$ 乘积测度。} \ $X\perp Y\iff P_{X,Y}=P_X\otimes P_Y$。

测度论有乘积测度，但没有“独立”这个\textbf{语义}：独立说的不只是“联合分布是乘积”，而是\textbf{“知道 $X$ 不改变对 $Y$ 的信念”}。这个信息层面在测度论里没有对应物。独立性派生的一切，测度论里一个都没有：

\begin{center}
\begin{tabular}{L{4cm}L{8cm}}
\toprule
\textbf{产物} & \textbf{内容} \\
\midrule
$0$-$1$ 律 & 尾 $\sigma$-代数上的事件概率只能 0 或 1 \\
Borel--Cantelli & $\sum P(A_n)<\infty\Rightarrow$ 只有有限个发生（a.s.） \\
Kolmogorov 扩展 & 相容的有限维分布 $\Rightarrow$ 存在无穷维过程 \\
de Finetti & 可交换序列 $=$ iid 的混合 \\
\bottomrule
\end{tabular}
\end{center}

\subsection{BR4.3 第 1 层：泛函分析}

\subsubsection{$\mathbb{E}$ 是线性泛函，特殊在“正 $+$ 保 1”}

$\mathbb{E}:L^1(P)\to\mathbb{R}$ 是\textbf{正}的（$X\ge0\Rightarrow\mathbb{E}X\ge0$）、\textbf{保常数}（$\mathbb{E}1=1$）、\textbf{范数 1}。Riesz 表示给 $C(K)^*=\mathcal{M}(K)$，概率测度正是其中“正 $+$ 总质量 1”的那些。\textbf{“期望是线性泛函”的信息量全在“正 $+$ 保 1”，不在“线性”。}

\subsubsection{条件期望 $=$ 投影（$L^2$ 专属）}

$\mathbb{E}[\cdot\mid\mathcal{G}]:L^2(\mathcal{F})\to L^2(\mathcal{G})$ 是\textbf{正交投影}：
\[
\lVert X-\mathbb{E}[X\mid\mathcal{G}]\rVert_2=\min_{Y\in L^2(\mathcal{G})}\lVert X-Y\rVert_2 .
\]
证明即正交性：对任意 $Y\in L^2(\mathcal{G})$，
\[
\mathbb{E}\big[(X-\mathbb{E}[X\mid\mathcal{G}])Y\big]
=\mathbb{E}[XY]-\mathbb{E}\big[\mathbb{E}[X\mid\mathcal{G}]Y\big]=0 .
\]

\textbf{价值在于搬运}：勾股定理、最小距离、正交分解全部搬进概率论。条件期望不是“一个定义”，它有一个几何图像——\textbf{在子空间上垂直投影}。

\textbf{副作用（全方差公式就是勾股定理）}：
\[
\mathrm{Var}(X)=\mathrm{Var}\big(\mathbb{E}[X\mid\mathcal{G}]\big)+\mathbb{E}\big[\mathrm{Var}(X\mid\mathcal{G})\big].
\]

\textbf{必须打的补丁}：在 $L^1$ 上它只是\textbf{压缩}（$\lVert\mathbb{E}[X\mid\mathcal{G}]\rVert_1\le\lVert X\rVert_1$），\textbf{不是投影}——$L^1$ 不严格凸，最近点不唯一。\textbf{“投影”这个几何图像 $L^2$ 专属。}

\subsubsection{马尔可夫过程 $=$ 压缩半群 $+$ 生成元}

$(P_tf)(x)=\mathbb{E}[f(X_t)\mid X_0=x]$，满足 $P_{t+s}=P_tP_s$、$P_0=I$、\textbf{保常数}、\textbf{保正}。

生成元 $Lf=\lim_{t\to0}\frac{P_tf-f}{t}$，形式地 $P_t=e^{tL}$。\textbf{Hille--Yosida}：压缩半群 $\iff$ 生成元（闭、稠密定义域、耗散）。

\textbf{这与 BR2/BR3 的“$D$ 无界 $+$ 自伴扩张”是同一件事的概率版}：$L$ 也无界，定义域就是“这个函数能不能当初始条件”；而“选边界条件”在这里变成“选过程的边界行为（吸收/反射）”。

\subsubsection{弱-$*$ 收敛与胎紧（一个必须纠正的术语）}

\[
\mu_n\Rightarrow\mu\iff\forall f\in C_b:\ \int f\,d\mu_n\to\int f\,d\mu .
\]
\textbf{这是 $C_b^*$ 中的 weak-$*$ 收敛，不是 $L^1$ 的弱收敛。} 泛函分析里真正的弱收敛是 $\sigma(L^1,L^\infty)$。两个都叫 weak，但一个是 $C_b$ 的对偶对，一个是 $L^\infty$ 的对偶对。\textbf{混了会在“用 $C_b$ 还是 $C_0$”上出错：概率论一律用 $C_b$。}

\textbf{紧性靠 Prohorov}：$\{\mu_n\}$ 弱-$*$ 相对紧 $\iff$ \textbf{胎紧}（tight）：对任意 $\epsilon>0$ 存在紧集 $K$ 使 $\mu_n(K)\ge1-\epsilon$ 对所有 $n$ 成立。

\textbf{“质量不跑掉”就是胎紧}。这是概率论对泛函分析最重要的贡献之一：把抽象紧性变成\textbf{可检验条件}。典型反例：$\delta_n$ 在 $\mathbb{R}$ 上跑向无穷，不胎紧、不收敛。

\subsubsection{鞅：条件期望的动力学}

\[
M_t\ \text{是鞅}\iff\mathbb{E}[M_{t+s}\mid\mathcal{F}_t]=M_t .
\]
\textbf{它是“条件期望的动力学”}：上一节说条件期望是投影（静态），鞅说这个投影随时间不动。\textbf{鞅 $=$ 投影沿时间轴的轨道。}

它不是“没变化”，是\textbf{“没有可预测的漂移”}——已知全部信息下，下一步的期望等于现在。这是概率论自己的“导数”概念。

\begin{center}
\begin{tabular}{L{3.4cm}L{9cm}}
\toprule
\textbf{产物} & \textbf{内容} \\
\midrule
Doob 分解 & 下鞅 $=$ 鞅 $+$ 可预测增过程（趋势 $+$ 噪声的典范分解） \\
可选停时 & 停时处取值的期望仍等于初值（不能靠停时赢公平游戏） \\
鞅收敛 & $L^1$ 有界鞅 a.s. 收敛 \\
$H^1$-$BMO$ 对偶 & Fefferman 定理 \\
鞅变换 & $\sum H_k\Delta M_k$ 仍是鞅 \\
\bottomrule
\end{tabular}
\end{center}

\subsection{BR4.4 第 2 层：调和分析}

\subsubsection{特征函数 $=$ 傅里叶变换}

$\varphi_X(t)=\mathbb{E}[e^{itX}]=\hat{P_X}(t)$。三个推论按顺序：

\textbf{Bochner 定理}（注意：是“正定函数”的 Bochner，不是“Bochner 积分”）：连续正定 $\varphi$（$\varphi(0)=1$，$\sum c_i\bar c_j\varphi(t_i-t_j)\ge0$）$\iff$ $\varphi$ 是某概率测度的特征函数。\textbf{这是“什么样的函数能当特征函数”的完整回答。}

\textbf{L\'evy 连续性}：$\varphi_n\to\varphi$ 逐点（$\varphi$ 在 $0$ 连续）$\Rightarrow\mu_n\Rightarrow\mu$。这是 CLT 的收尾工具。

\textbf{L\'evy--Khintchine}（无穷可分分类）：
\[
\log\varphi(t)=iat-\frac{\sigma^2t^2}{2}+\int_{\mathbb{R}}\left(e^{itx}-1-itx\mathbf{1}_{|x|<1}\right)\nu(dx).
\]
L\'evy 三元组 $(a,\sigma^2,\nu)$ 完全刻画。\textbf{这是“独立增量过程的分类定理”}：高斯（$\nu=0$）$+$ 泊松（$\sigma=0$，$\nu$ 有限）$+$ 稳定（$\nu$ 幂律）全在这一个公式里，也是 BR2.5 稳定分布的完整骨架。

\subsubsection{鞅变换 $\leftrightarrow$ 奇异积分（唯一“定理级”的双向对应）}

\textbf{鞅变换}：$M$ 鞅，$H$ 可预测（$H_k\in\mathcal{F}_{k-1}$），
\[
(H\cdot M)_n=\sum_{k\le n}H_k\,(M_k-M_{k-1}).
\]
它仍是鞅，因为
\[
\mathbb{E}[H_k\Delta M_k\mid\mathcal{F}_{k-1}]
=H_k\,\mathbb{E}[\Delta M_k\mid\mathcal{F}_{k-1}]=0 .
\]
\textbf{“可预测性”是全部关键}：$H_k$ 必须在下注之前就知道，否则“鞅变换还是鞅”就崩——这是“你不能靠内幕交易赢公平游戏”的数学版。

\textbf{Burkholder 不等式}（$1<p<\infty$）：
\[
c_p\lVert H\rVert_{L^p}\le\lVert H\cdot M\rVert_p\le C_p\lVert H\rVert_p .
\]
它对应 \textbf{Hilbert 变换}
\[
Hf(x)=\mathrm{p.v.}\int\frac{f(y)}{x-y}dy
\]
的 $L^p$ 有界性。具体对应：把 $\mathbb{R}$ 上的 dyadic 区间族看成鞅的滤子，Hilbert 变换的截断版本就是一个鞅变换。

\textbf{所以 Burkholder 不等式 $=$ 概率版的 Calder\'on--Zygmund 奇异积分理论。} 这是“双向地图”最实在的一个实例——不是类比，是定理级的等价：$L^p$ 有界可以走概率证明，鞅的 $L^p$ 估计可以走分析插值。

\textbf{顺带}：Fefferman 对偶 $(H^1)^*=BMO$。概率的“原子”（$L^2$ 有界、条件期望为零的鞅差）与分析的“原子”（支撑在区间、均值为零的函数）是同一个东西的两套坐标。John--Nirenberg 不等式（$BMO$ 的指数可积性）也是从鞅的好 $\lambda$ 不等式来的。

\subsection{BR4.5 第 3 层：几何}

\subsubsection{协方差 $=$ 度量}

扩散 $dX_t=b\,dt+\sigma\,dW_t$ 的生成元
\[
L=\sum_ib_i\partial_i+\frac12\sum_{i,j}(\sigma\sigma^T)_{ij}\partial_i\partial_j .
\]
当 $b=0$、$\sigma\sigma^T=2g^{-1}$ 时 $L=\Delta_g$。故
\[
a:=\sigma\sigma^T\ \longleftrightarrow\ g^{-1}
\qquad(\text{协方差矩阵}\leftrightarrow\text{逆度量}).
\]
上一轮 $\Delta_g$ 坐标式里的 $\sqrt{|g|}$ 对应\textbf{平稳测度的密度}——这解释了为什么流形上布朗运动的平稳分布\textbf{不是 Lebesgue 而是体积测度}。

\subsubsection{热核 $=$ 转移密度（接回 BR1）}

$p_t(x,y)$ 是 $e^{t\Delta_g}$ 的积分核，短时渐近（Varadhan）
\[
p_t(x,y)\sim\frac{e^{-d(x,y)^2/4t}}{(4\pi t)^{n/2}}\qquad(t\to0).
\]
\textbf{短时间内的转移概率由测地距离决定}。而 BR1 算的 $\arccos\langle p,q\rangle$ 就是 $S^n$ 上的 $d(p,q)$——\textbf{大圆 $=$ 热核衰减最快的路径 $=$ 最可能被扩散走的路}。几何、概率、分析在这里是同一句话。

\subsubsection{曲率 $\to$ 集中：中间夹着 $\Gamma$ 与 $\Gamma_2$}

这是“Ric $\to$ 集中”的真实通道。\textbf{Carr\'e du champ}：
\[
\Gamma(f,g)=\frac12\big(L(fg)-fLg-gLf\big),
\]
对 $L=\Delta_g$ 就是 $\Gamma(f,g)=\langle\nabla f,\nabla g\rangle_g$——\textbf{它就是 BR3 的度量}。二阶算子
\[
\Gamma_2(f,f)=\frac12\big(L\Gamma(f,f)-2\Gamma(f,Lf)\big).
\]
\textbf{Bakry--\'Emery 曲率条件} $\Gamma_2(f,f)\ge\kappa\Gamma(f,f)$：

\begin{itemize}
  \item $\kappa>0$ $\Rightarrow$ Bonnet--Myers 型直径界 $+$ L\'evy 型测度集中 $+$ Poincar\'e 谱隙；
  \item $L=\Delta_g$、体积测度：$\kappa=\mathrm{Ric}$ 的下界；
  \item 测度 $e^{-V}dx$（$L=\Delta_g-\nabla V\cdot\nabla$）：$\kappa=\mathrm{Ric}+\mathrm{Hess}V$ 的下界——Bakry--\'Emery 张量。
\end{itemize}

\textbf{真正的等式链是}
\[
\mathrm{Ric}\ (\text{或 }\mathrm{Ric}+\mathrm{Hess}V)\ \ge\ \kappa
\iff\Gamma_2\ge\kappa\Gamma
\Longrightarrow\ \text{集中/直径界/谱隙}.
\]
\textbf{“$\kappa>0\Rightarrow$ 集中”的直觉}：正 $\Gamma_2$ 压住“梯度的梯度”，函数不能又陡又持久，于是 Lipschitz 函数的值紧贴均值。这是测度集中的微分版本。

\subsubsection{Wiener 测度：两个反直觉事实}

$C([0,1])$ 上的 Wiener 测度 $\mathbb{W}$，特征函数
\[
\mathbb{E}_{\mathbb{W}}\left[e^{i\langle f,\omega\rangle}\right]=e^{-\frac12\lVert f\rVert_H^2}
\]
（$H$ 是 Cameron--Martin 空间）——BR2“高斯自对偶”的无穷维版本。

\textbf{（a）Cameron--Martin}：$\mathbb{W}$ 与平移 $\mathbb{W}_h$ 等价（RN 导数存在）$\iff h\in H$；$h\notin H$ 时\textbf{互相奇异}。

\textbf{（b）更反直觉}：$\mathbb{W}(H)=0$。\textbf{“可微方向”本身是零测集}，典型样本不在任何“好方向”上。

这解释了“无穷维不是有限维取极限”：无穷维\textbf{不存在 Lebesgue 测度}，Wiener 测度只在一个稠密零测的子空间 $H$ 上准不变。

\textbf{Malliavin 微积分}是 BR2/BR3 那个 $D=d/dx$ 的无穷维版：$D_hF=\lim_{\epsilon\to0}\frac{F(\omega+\epsilon h)-F(\omega)}{\epsilon}$，但只有 $h\in H$ 的平移才有意义。\textbf{它就是“在一个几乎所有方向都奇异的空间上做微分”的技术}，产物是分部积分（Skorokhod 积分）与 H\"ormander 定理。

\subsection{BR4.6 第 4 层：PDE}

\subsubsection{Feynman--Kac}

\[
\partial_tu=Lu-Vu,\ u(0,x)=f(x)
\quad\Longleftrightarrow\quad
u(t,x)=\mathbb{E}_x\left[f(X_t)\exp\left(-\int_0^tV(X_s)ds\right)\right].
\]
三种读法：
\begin{enumerate}
  \item \textbf{PDE 的概率表示}：解 $=$ 期望，两边定性可互换。
  \item \textbf{半群的扰动} $e^{t(L-V)}$：$V$ 是杀死率，$e^{-\int V}$ 是存活概率。
  \item \textbf{Trotter 公式}：$e^{t(L-V)}=\lim_n\left(e^{tL/n}e^{-tV/n}\right)^n$——“走一步 $+$ 掷骰子决定死不死”。
\end{enumerate}
这把 $e^{t\cdot}$ 一路焊死：$e^{tA}$（矩阵）$\to e^{tL}$（算子半群）$\to e^{t(L-V)}$（Feynman--Kac）。

\subsubsection{大偏差：枢纽是 Legendre 变换}

\[
I(x)=\sup_{\theta\in\mathbb{R}}\big(\theta x-\log M(\theta)\big),\qquad M(\theta)=\mathbb{E}[e^{\theta X}].
\]
\textbf{这就是 Legendre 变换（凸共轭）}，$I=M^*$。准确链条是
\[
\log M\ \xrightarrow{\text{Legendre}}\ I\ \xrightarrow{\text{路径版}}\ 
\frac12\int\lVert\dot\phi-b(\phi)\rVert^2 .
\]
\textbf{Cramér}（iid）：$-\frac1n\log P(\bar X_n\in A)\to\inf_{x\in A}I(x)$。

\textbf{为什么 Legendre 是枢纽}：
\begin{itemize}
  \item 它是凸分析的对偶（与 BR2“选基对角化”同类：\textbf{选对偶坐标}）；
  \item $I(x)=\sup_\theta(\theta x-\log M)$ 就是“用 $\theta$ 当 Lagrange 乘子、在 $\mathbb{E}X=x$ 下最大化熵”——\textbf{速率函数 $=$ 约束下的熵减}，这条把大偏差与最大熵接起来。
\end{itemize}

\textbf{Freidlin--Wentzell}（小噪声 $dX=b\,dt+\sqrt\epsilon\,dW$）：$I(\phi)=\frac12\int_0^T\lVert\dot\phi_s-b(\phi_s)\rVert^2ds$。\textbf{这就是作用量}——最小作用量原理变成“最可能的偏离路径”。

\textbf{粘性解}：Hamilton--Jacobi $-\partial_tu+H(x,\nabla u)=0$ 的解可由 $-\epsilon\log$（大偏差）逼近。\textbf{粘性解 $=$ 大偏差极限的 PDE}，链条是“大偏差 $\to$ 变分 $\to$ HJ $\to$ 粘性解”。

\subsubsection{SPDE（只标注边界，不深入）}

随机热方程 $\partial_tu=\Delta u+\dot W$ 解是高斯过程，好办。乘性噪声 $\partial_tu=\Delta u+\sigma(u)\dot W$ 在 $d\ge2$ 时 \textbf{$\sigma(u)\dot W$ 无定义}（$u$ 不够光滑，两个分布不能相乘）。Hairer 的正则结构（2014）给“病态乘积”造代数框架，把 $u$ 提升到正则结构里让乘积有意义；KPZ $\partial_th=\partial_x^2h+(\partial_xh)^2+\xi$ 是原型。

\textbf{入门者只需知道}：它解决的是“乘积无定义”这个纯代数问题，与 BR3 的“$d^2=0$ 需要 de Rham 复形”同类（给不够好的对象找个够好的容器）。

\subsection{BR4.7 第 5 层：代数、数论、组合}

\subsubsection{随机矩阵}

\textbf{Wigner 半圆律}：$N\times N$ 对称随机矩阵（上三角 iid，方差 $1/N$），经验谱分布
\[
\mu_N=\frac1N\sum_{i=1}^N\delta_{\lambda_i}\ \longrightarrow\
\sigma(x)=\frac{1}{2\pi}\sqrt{4-x^2}\,\mathbf{1}_{|x|\le2}.
\]
\textbf{为什么半圆}：矩 $\int x^{2k}d\sigma=C_k$（Catalan 数），Catalan 数数的是\textbf{非交叉配对}，而非交叉配对正是自由概率的自由累积量。\textbf{矩方法的胜利}：分布由矩唯一决定（Carleman 条件），矩由组合计数给出。

\textbf{Tracy--Widom}：$N^{2/3}(\lambda_{\max}-2)\Rightarrow\mathrm{TW}_2$，是 Painlev\'e II 的解。\textbf{注意尺度}：平均是 $N^{1/2}$ 窗口，极大值涨落是 $N^{-1/6}$ 窗口（乘 $N^{2/3}$ 才稳）。这个 $1/6$ 是 KPZ 普适类的指纹。

\subsubsection{自由概率：第三种独立性}

\textbf{自由独立性}：$A,B$ 自由独立 $\iff$ 对任意无常数项的非交换多项式 $P,Q$，$\mathbb{E}[P(A)Q(B)]=0$。

这个定义看起来怪，它的作用是把经典 CLT 的机制换掉：

\begin{center}
\begin{tabular}{L{4cm}L{8cm}}
\toprule
\textbf{独立性} & \textbf{CLT 极限} \\
\midrule
经典独立 & 高斯 \\
自由独立 & \textbf{半圆律} \\
Boolean 独立 & Bernoulli 律 \\
单调独立 & arcsine 律 \\
\bottomrule
\end{tabular}
\end{center}

\textbf{“CLT 的极限分布完全由‘什么算独立’决定”} —— 这是“概率论 $\to$ 代数”反馈的核心。而它反过来解释几何：半圆律是随机矩阵的极限谱，所以大随机矩阵 $=$ 自由独立的算子，这把算子代数与概率钉死。

\subsubsection{概率方法与概率数论}

\textbf{概率方法（Erd\H{o}s）}：Ramsey 下界 $R(k,k)>2^{k/2}$——随机 2-着色，任一 $K_k$ 单色的概率 $<1$，故\textbf{存在}好着色。\textbf{它暴露一个哲学点：概率论能给“存在”，不能给“构造”}——非构造存在性。

\textbf{Erd\H{o}s--Kac}：
\[
\frac{\omega(n)-\log\log n}{\sqrt{\log\log n}}\Rightarrow N(0,1)
\]
（$\omega(n)$ 是 $n$ 的不同素因子个数）。不同素数的可除性近似独立，启发式的独立性直接给出正确答案，严格化靠筛法。\textbf{“随机化”不是近似工具，是发现工具。}

\subsection{BR4.8 六个反馈产物（补全清单）}

\subsubsection{（一）五种收敛——地图最容易漏的一条}

\begin{center}
\begin{tabular}{L{3.2cm}L{5.5cm}L{3.5cm}}
\toprule
\textbf{收敛} & \textbf{定义} & \textbf{强度} \\
\midrule
几乎必然 a.s. & $P(X_n\to X)=1$ & 强 \\
依概率 & $\forall\epsilon:P(|X_n-X|>\epsilon)\to0$ & 中 \\
依分布 $\Rightarrow$ & $\mathbb{E}[f(X_n)]\to\mathbb{E}[f(X)]$ & 弱 \\
$L^p$ & $\mathbb{E}|X_n-X|^p\to0$ & 中（依 $p$） \\
完全收敛 & $\sum_nP(|X_n-X|>\epsilon)<\infty$ & 强 \\
\bottomrule
\end{tabular}
\end{center}

关系：a.s. $\Rightarrow$ 依概率 $\Rightarrow$ 依分布；$L^p\Rightarrow$ 依概率；完全收敛 $\Rightarrow$ a.s.；依概率 $\Rightarrow$ 存在子列 a.s.。\textbf{反向一般不成立。}

\textbf{为什么这是“反馈回数学”}：测度论原本只有“测度收敛”一个概念。概率论造出五种“随机变量收敛”，每种对应一个场景（大数定律用 a.s.、CLT 用依分布、鞅用 a.s.、$L^p$ 估计用 $L^p$）。\textbf{这是随机化逼出的概念分层，现在已是测度论的标准内容。}

\subsubsection{（二）耦合 $\to$ 最优传输（概率论自己生的）}

\textbf{耦合}：$\Pi(\mu,\nu)=\{\pi:\text{边缘分别是 }\mu,\nu\}$。
\[
W_p(\mu,\nu)=\inf_{\pi\in\Pi(\mu,\nu)}\left(\int d(x,y)^p\,d\pi\right)^{1/p}.
\]
\textbf{$W_p$ 就是“最优耦合”}——最优传输从概率这边看，就是耦合的定量化。\textbf{这是“概率 $\to$ 几何”最纯的例子}，不是几何强加的。

\textbf{Brenier 定理}：$\mu\ll\lambda$ 时 $W_2$ 最优耦合是确定性映射 $T=\nabla\phi$（$\phi$ 凸）——\textbf{最优传输 $=$ 凸函数的梯度}。

\textbf{Lott--Sturm--Villani}：在 $W_2$ 下，$\mathrm{Ric}\ge K\iff$ 熵沿 $W_2$ 测地线 $K$-凸。\textbf{“曲率 $=$ 熵凸性”的精确版本}，而驱动它的正是“耦合”这个纯概率概念。

\subsubsection{（三）Dirichlet 空间 $=$ 概率版 Sobolev}

布朗运动与调和函数：
\[
u(x)=\mathbb{E}_x[f(X_{\tau_D})]\quad\text{解 Dirichlet 问题 }(\Delta u=0,\ u|_{\partial D}=f),
\]
$\tau_D$ 是出 $D$ 的首中时。\textbf{Dirichlet 空间}（$L^2$ 上对称 Markov 半群 $+$ 能量形式 $\mathcal{E}(f,g)=\langle f,-Lg\rangle$）\textbf{就是“概率版的 Sobolev 空间”}：BR3 的 Sobolev 是“让 $D$ 有界”，Dirichlet 空间是“让 $L$ 自伴”——同一件事的两个名字。

\subsubsection{（四）(五)(六) 速记}

\begin{itemize}
  \item \textbf{乘子理论}：独立和的极限 $\leftrightarrow$ 傅里叶乘子；Marcinkiewicz / H\"ormander--Mikhlin。概率给出“为什么乘子有界”的概率证明。
  \item \textbf{随机矩阵 $\to$ 因子}：自由概率让 von Neumann 因子有了概率语言（Connes 嵌入问题）。
  \item \textbf{测度集中}：高维测度的几何现象，反馈成度量测度空间与 RCD 空间理论。
\end{itemize}

\subsection{BR4.9 一句话收束}

\textbf{数学 $\to$ 概率}：测度变 $P$，积分变 $\mathbb{E}$，\textbf{正交投影变条件期望}（$L^2$），\textbf{压缩半群变马尔可夫过程}，\textbf{度量 $g^{-1}$ 变协方差 $\sigma\sigma^T$}，\textbf{Laplace--Beltrami 变生成元}，\textbf{$\Gamma_2$ 变集中性}，\textbf{Legendre 变换变速率函数}。

\textbf{概率 $\to$ 数学}：RN 导数、五种收敛、胎紧性、条件期望作为投影、鞅空间与 $H^1$-$BMO$、Burkholder $\leftrightarrow$ Calder\'on--Zygmund、Wiener 测度与无穷维高斯、耦合 $\to W_2\to$ Lott--Sturm--Villani、Feynman--Kac、Legendre $\to$ 粘性解、正则结构、随机矩阵与自由概率、概率方法。

\textbf{而整份地图的经线只有三根}：

\begin{center}
\begin{tabular}{L{2.4cm}L{3.2cm}L{7cm}}
\toprule
\textbf{原语} & \textbf{代数身份} & \textbf{派生物} \\
\midrule
独立性 & 乘积测度 & $0$-$1$ 律、Borel--Cantelli、CLT \\
条件化 & $L^2$ 正交投影 & 条件期望、Markov 性、Bayes \\
鞅 & 条件期望的动力学 & Doob 分解、停时、$H^1$-$BMO$ \\
\bottomrule
\end{tabular}
\end{center}

\textbf{测度论一个都没有。} 概率论不是测度论的应用，它是“在测度论上装了时间轴与信息轴”之后长出来的新学科。

\end{document}